\section{Preliminary}\label{sec:preliminary}

\begin{figure}[t]
\centering
\footnotesize
\setlength{\jot}{0pt}

\begin{minipage}[t]{0.4\columnwidth}
\[
\begin{alignedat}{2}
i &\;::=\; \textit{literal}; &\quad& \text{Identifier} \\
il &\;::=\; \{i\}; & & \text{List of identifiers} \\
pred_{il} &\;::=\; predicate; & & \text{Predicate such that } \\
 & & & \quad FV(pred_{il}) \subseteq il \\
M &\;::=\; {Context}, & & \text{Formal model} \\
 & \quad {Machine}; & & \\
 & & & \\
Context &\;::=\; i, & & \text{Context identifier} \\
 & \quad [\text{``extends''},i], & & \text{Extended context} \\
 & \quad \text{``sets''},s, & & s ::= il \\
 & \quad \text{``constants''},c, & & c ::= il \\
 & \quad \text{``axioms''},A, & & A ::= \{pred_{s \cup c}\} \\
 & \quad \text{``theorems''},T; & & T ::= \{pred_{s \cup c}\} 
\end{alignedat}
\]
\end{minipage}
\hfill
\begin{minipage}[t]{0.57\columnwidth}
\[
\begin{alignedat}{2}
Machine &\;::=\; i, &\quad& \text{Machine identifier} \\
 & \quad [\text{``refines''},i], & & \text{Refined machine} \\
 & \quad [\text{``sees''},i], & & \text{Seen context} \\
 & \quad \text{``variables''},v, & & v ::= il\\
 & \quad \text{``invariants''},I, & & I ::= \{pred_{s \cup c \cup v}\} \\
 & \quad \text{``variants''},N, & & N::= \{pred_{s \cup c \cup v}\} \\
 & \quad \text{``theorems''},T, & & T ::= \{pred_{s \cup c \cup v}\} \\
 & \quad \text{``events''},\mathcal{E}; & & \mathcal{E} ::= \{Event\} \\
 & & & \\
Event &\;::=\; i, & & \text{Event identifier} \\
 & \quad \text{``any''},x, & & \text{Event parameters, } x ::= il \\
 & \quad \text{``where''},G, & & \text{Event guards, } G ::= \{pred_{s \cup c \cup v \cup x}\} \\
 & \quad \text{``then''},v \;:|\; BA; & & \text{Actions, } BA ::= \{pred_{s \cup c \cup v \cup x \cup v'}\} \\
\end{alignedat}
\]
\end{minipage}

\caption{The grammar of Event-B, where \(FV(pred)\) is the set of free variables in \(pred\).}
\label{fig:grammar}
\end{figure}


\textbf{The Event-B Notation.}
\autoref{fig:grammar} summarizes the grammar of Event-B \cite{abrial2010modeling}.
An Event-B model $M$ is a discrete transition system grounded in set theory and first-order logic, where all types are represented as sets.
For instance, $STRING$ denotes the set of all strings, and $A \rightarrow B$ the set of total functions with domain $A$ and range $B$.

A model $M$ consists of \emph{contexts} and \emph{machines}.
Let $il$ be a list of identifiers. Sets $s$, constants $c$, variables before and after an event $v$ and $v'$, and event parameters $x$ are defined over $il$ and are disjoined.
A predicate $pred$ with $FV(pred) = il$ is denoted by $pred_{il}$, where $FV(pred)$ is the set of free variables in $pred$.
A list of predicates using only variables in $il$ is denoted by $\{pred_{il}\}$, which can be used to express components like axioms, $A ::= \{pred_{s \cup c}\}$.

\noindent
\textbf{Correctness by Construction.}
Let $REQ_M$ denote the requirements of formal model $M$, and $PO_M$ the set of its proof obligations.
Each $req \in REQ_M$ is associated with a predicate $Covered(M, req) \in \{\mathsf{true}, \mathsf{false}\}$, which holds iff $req$ is covered in $M$ as an invariant, variant, axiom, or event.
Similarly, each $po \in PO_M$ has a predicate $Discharged(M,po) \in \{\mathsf{true}, \mathsf{false}\}$, which holds iff $po$ is discharged.
We define a predicate $Correct(M)$ to indicate that $M$ is \emph{correct by construction}, namely:
{\footnotesize\begin{equation}
Correct(M) \;\triangleq\;
\Bigl(\bigwedge_{req \in REQ_M} Covered(M, req)\Bigr)
\;\wedge\;
\Bigl(\bigwedge_{po \in PO_M} Discharged(M,po)\Bigr)
\label{eq:correctness}
\end{equation}}
